%-------------------------
% Resume in LaTeX
% Author: Zeep Jia
%-------------------------

\documentclass[a4paper,11pt]{article}

\usepackage[UTF8]{ctex}
\usepackage[left=0.6in,right=0.6in,top=0.5in,bottom=0.5in]{geometry}
\usepackage{titlesec}
\usepackage{enumitem}
\usepackage{hyperref}
\usepackage{fontawesome5}
\usepackage{xcolor}
\usepackage{tabularx}
\usepackage{tikz}

% 颜色定义
\definecolor{primary}{RGB}{31, 78, 121}
\definecolor{secondary}{RGB}{89, 89, 89}
\definecolor{accent}{RGB}{0, 102, 153}

% 超链接样式
\hypersetup{
    colorlinks=true,
    linkcolor=accent,
    urlcolor=accent,
    pdfborder={0 0 0}
}

% 段落设置
\setlength{\parindent}{0pt}
\setlength{\parskip}{0pt}

% 标题样式
\titleformat{\section}
  {\large\bfseries\color{primary}}
  {}
  {0em}
  {}
  [\vspace{-0.5em}\textcolor{primary}{\rule{\linewidth}{1pt}}]

\titlespacing*{\section}{0pt}{0.6em}{0.4em}

% 列表样式
\setlist[itemize]{leftmargin=1.2em, itemsep=0pt, parsep=0pt, topsep=2pt}

% 自定义命令
\newcommand{\entry}[4]{
  \textbf{#1} \hfill \textcolor{secondary}{\small #2} \\
  \textit{\small #3} \hfill \textcolor{secondary}{\small #4}
}

\newcommand{\projectentry}[3]{
  \textbf{#1} \hfill \textcolor{secondary}{\small #2} \\
  \textit{\small #3}
}

\begin{document}

%==== 头部信息 ====%
\begin{center}
    {\Huge\bfseries\color{primary} 贾泽澎 \,|\, Zeep Jia}
    
    \vspace{0.3em}
    {\large\color{secondary} 求职意向：产品经理实习 / AI 产品实习}
    
    \vspace{0.4em}
    \small
    \faPhone\ 158-9989-2439 \quad
    \faWeixin\ Jzp15899892439 \quad
    \faEnvelope\ \href{mailto:zepengjia611@gmail.com}{zepengjia611@gmail.com} \quad
    \faGithub\ \href{https://github.com/MOLIYUANQUAN}{MOLIYUANQUAN}
    
    \vspace{0.1em}
    \faMapMarker*\ 深圳 / 曼彻斯特
\end{center}

\vspace{-0.3em}

%==== 个人简介 ====%
\section{个人简介}
曼彻斯特大学计算机科学本科在读，关注 AI Agent、校园社区与效率工具。具备从 0 到 1 产品落地、用户增长与跨职能协作经验，曾发起校园社区小程序「海外留子圈」，组建 10+ 人团队，并在曼大中国留学生群体中积累近 \textbf{1000 名}真实用户。

%==== 教育背景 ====%
\section{教育背景}
\entry{曼彻斯特大学 The University of Manchester}{2024.09 -- 2027.06}{BSc Computer Science 理学学士，计算机科学}{大二在读}

\vspace{0.2em}
\small\textbf{核心课程：}数据结构与算法、软件工程、数据库系统、分布式系统、机器学习、基于知识的 AI、数据科学、系统架构

%==== 产品与项目经历 ====%
\section{产品与项目经历}

\projectentry{海外留子圈｜校园 C2C 与社交平台}{2025.09 -- 至今}{创始人 \& 产品经理}
\begin{itemize}
    \item 从 0 到 1 发起面向海外中国留学生的校园社区微信小程序，聚焦二手交易、校园信息流、轻社交与活动组局
    \item 组建并管理 10+ 人学生团队，负责产品方向、MVP 边界、版本节奏与跨职能协作，推动产品于 2026.03 上线
    \item 主导核心功能设计：二手帖发布/浏览、Feed 流、评论互动、私信会话、组局活动发布与报名等模块
    \item 基于用户反馈持续优化产品体验，通过微信群与校园社群完成冷启动，积累近 \textbf{1000 名}真实用户
\end{itemize}

\vspace{0.3em}
\projectentry{ReproShip｜论文代码复现 Agent}{2026.03}{产品负责人}
\begin{itemize}
    \item 针对论文代码复现中的环境配置复杂、依赖冲突和报错排查成本高等问题，提出代码复现 Agent 产品方向
    \item 设计「拉取代码 → 沙箱运行 → 日志诊断 → 自动修复建议 → 复现报告」核心任务流，推进可演示原型落地
\end{itemize}

\vspace{0.3em}
\projectentry{早期产品与技术原型探索}{2025 -- 2026}{产品探索 / 开发实践}
\begin{itemize}
    \item \textbf{AIHub}：基于 Expo / React Native 搭建多模型 AI 工作台移动端 MVP，探索统一入口与提示词模板
    \item \textbf{Arcflow}：AI 工作流工具，将写作、调研、开发等任务拆成清晰步骤，帮助用户跟着流程完成产出
    \item \textbf{EventLite}：基于 Java Spring Boot 开发校园活动管理平台，参与后端接口与数据模型设计
\end{itemize}

%==== 技能与工具 ====%
\section{技能与工具}
\begin{tabularx}{\textwidth}{@{}l X@{}}
\textbf{产品能力} & 需求分析、用户故事、PRD、MVP 定义、信息架构、竞品分析、用户反馈分析、版本规划、跨职能协作 \\[0.3em]
\textbf{增长运营} & 冷启动、校园社群运营、用户转化、内容供给、微信群传播、增长数据分析 \\[0.3em]
\textbf{AI 工具} & ChatGPT、Claude、Cursor、Claude Code、Codex、AI Agent Workflow \\[0.3em]
\textbf{技术理解} & Python、JavaScript/TypeScript、Java、SQL、React Native、微信小程序、Spring Boot、Node.js \\
\end{tabularx}

\end{document}
